\chapter{Chapter 2: Functions, definitions, and structures}\label{sec:02-functions-and-structures:00-index:1}

Functions \textbf{curry}: \texttt{add : Nat → Nat → Nat} is really \texttt{Nat → (Nat → Nat)}, a function that returns another function. So a ``two-argument function'' is just a one-argument function whose result is itself another function. It takes one argument at a time. (For readers already thinking categorically: this is the type-theoretic form of the Hom-set isomorphism \(\mathrm{Hom}(A\times B, C)\cong\mathrm{Hom}(A,\mathrm{Hom}(B,C))\). A two-argument map is the same data as a one-argument map into a space of maps. This is not needed to follow along --- the plain statement above is enough.) The interesting part of this chapter is \texttt{structure}, which is how algebraic data will be packaged.

\textbf{Learning objectives.} By the end of this chapter, bundle data (and proofs) into a \texttt{structure}, use field projection and the anonymous constructor, parameterize a \texttt{structure} by a type, and extend one \texttt{structure} with another via \texttt{extends}.

\section{Sections}\label{sec:02-functions-and-structures:00-index:2}

\begin{enumerate}
\def\labelenumi{\arabic{enumi}.}
\tightlist
\item
  \hyperref[sec:02-functions-and-structures:01-structure-basics:1]{\texttt{structure}: bundling data together}
\item
  \hyperref[sec:02-functions-and-structures:02-type-parameters:1]{Structures with type parameters}
\item
  \hyperref[sec:02-functions-and-structures:03-extending-structures:1]{Extending structures}
\end{enumerate}

% Auto-appended: input every other section in order.
\section{\texorpdfstring{\texttt{structure}: bundling data together}{structure: bundling data together}}\label{sec:02-functions-and-structures:01-structure-basics:1}

Chapter 1 fixed what a single term and a single function look like, down to the calculus of constructions itself. But a function of several arguments, as the chapter introduction just noted, is really a chain of one-argument functions --- \texttt{add : Nat → Nat → Nat} never actually receives two numbers at once. That is fine for arguments that stay independent, but plenty of the mathematics this book will build --- a point in the plane, later a group's carrier and operation and axioms together --- needs several pieces of data to travel \emph{together}, as one value, not as separate arguments a caller has to remember to keep in sync. \texttt{structure} is Lean's answer to that need.

A \texttt{structure} groups several pieces of data under one name. We will use this Lean feature constantly once we define groups and rings.

\begin{lstlisting}
structure Point where
  x : Nat
  y : Nat

def origin : Point := { x := 0, y := 0 }

#eval origin.x        -- 0
\end{lstlisting}

Key points:

\begin{itemize}
\tightlist
\item
  \texttt{Point.mk} is the automatically generated \textbf{constructor}. This is the function that actually builds a \texttt{Point} out of an \texttt{x} and a \texttt{y}. \texttt{\{ x := ..., y := ... \}} is shorthand for \texttt{Point.mk ... ...}. It names each field so the order cannot be mixed up.
\item
  There is an even shorter way to write this: \texttt{⟨0, 0⟩} builds the exact same \texttt{Point}. It simply lists the values in field-declaration order instead of naming them. Read \texttt{⟨\_\allowbreak{}, \_\allowbreak{}⟩} as \textbf{``here are the pieces, in order --- the constructor is to be inferred from context.''} Lean can always determine it, because the \emph{expected type} (here, \texttt{Point}, from \texttt{def origin : Point := ...}) indicates exactly which constructor and which fields must be filled in. This is also where the official name comes from: it is called the \textbf{anonymous constructor}, because \texttt{Point.mk} is never written explicitly --- it remains anonymous, and Lean infers it from context. \texttt{⟨\_\allowbreak{}, \_\allowbreak{}⟩} recurs constantly from Chapter 3 onward, for proofs as much as for data.
\item
  \texttt{p.x} is \textbf{field projection} notation, shorthand for \texttt{Point.x p}. \texttt{origin.x} above is exactly this projection applied to \texttt{origin}.
\end{itemize}

\begin{lstlisting}
def shift (p : Point) (dx dy : Nat) : Point :=
  { x := p.x + dx, y := p.y + dy }

#eval (shift origin 3 4).y   -- 4
\end{lstlisting}

\texttt{shift} shows a structure used on \emph{both} sides of a function: it takes a \texttt{Point} in (reading its fields back out with the same \texttt{p.x}/\texttt{p.y} projection notation) and builds a new one via \texttt{\{ x := ..., y := ... \}}, the same field-naming syntax \texttt{origin} used above.

Note also that structures can bundle \emph{proofs} alongside data, not just data. This is exactly how a group will be defined later --- a carrier type, an operation, and proofs that the operation satisfies the group axioms, all in one \texttt{structure}.

\begin{mathreading}

\texttt{structure Point where x : Nat; y : Nat} is the Cartesian product \(\mathrm{Point} = \mathbb{N} \times \mathbb{N}\), with \texttt{x} and \texttt{y} playing the role of the two projections \(\pi_1, \pi_2 : \mathrm{Point} \to \mathbb{N}\). Here \texttt{p.x} computes \(\pi_1(p)\), \texttt{p.y} computes \(\pi_2(p)\), and \texttt{\{ x := ..., y := ... \}} builds an element via the universal property of the product (a pair of maps into the factors determines a unique map into the product). More generally, a \texttt{structure} with fields of types \(A_1, \ldots, A_n\) is the \(n\)-fold product \(A_1 \times \cdots \times A_n\). Once fields are allowed to be \emph{proofs} (propositions), the same construction becomes a \textbf{subset cut out by conditions}: a structure \texttt{\{ data : D, proof : P data \}} is the dependent pair (subset) \(\{\, d \in D \mid P(d) \,\}\), categorically a subobject of \(D\).

\end{mathreading}

Here is the universal property itself, in this box's own notation: the unique \(h\) making both triangles commute is exactly what \texttt{\{ x := ..., y := ... \}} constructs when \(f\) and \(g\) are \(p\)'s two argument expressions.

\begin{center}
% Universal property of the categorical product A x B, in this box's own
% notation: C -> A (f), C -> B (g), and a unique mediating h : C -> A x B
% making both triangles commute via the projections p_A, p_B.
% Source: 02-functions-and-structures/01-structure-basics.md, "Categorical
% product (universal property)" Recall entry / Mathematical reading box.
\begin{tikzcd}
  & C \arrow[dl, "f"'] \arrow[dr, "g"] \arrow[d, dashed, "\exists ! h"] & \\
  A & A \times B \arrow[l, "p_A"] \arrow[r, "p_B"'] & B
\end{tikzcd}

\end{center}

\subsection{Sources, quoted}\label{sec:02-functions-and-structures:01-structure-basics:2}

Formal definitions and citations for this section, gathered here for reference (full entries in the \hyperref[sec:root:bibliography:1]{Bibliography}):

\begin{itemize}
\tightlist
\item
  \textbf{Record / tuple.} ``The simplest of these is pairs, or more generally tuples, of values \ldots{} the generalization from n-ary tuples to labeled records is equally straightforward'' (\cite{Pierce2002}, §11.6--§11.8). Picture it like this: an unlabeled tuple is a box packed in a fixed order --- you just have to remember ``first slot is the name, second is the age.'' A labeled record is the same box with every slot tagged directly: ``Name: \textbf{\emph{, Age: }}.'' Lean's \texttt{structure} is exactly this tagged version.
\item
  \textbf{Categorical product (universal property).} For a category with a product \(A \times B\) (with projections \(p_A, p_B\)) and any object \(C\) with morphisms \(f : C \to A\), \(g : C \to B\): ``there is exactly one morphism \(h : C \to A \times B\) such that \(f = p_A h\) and \(g = p_B h\)'' (\cite{Pareigis1970}, §1.11, p.~30).
\item
  Lean 4 documentation (\cite{LeanDocs}) --- the constructor/projection/anonymous-constructor mechanics described above.
\end{itemize}

\section{Structures with type parameters}\label{sec:02-functions-and-structures:02-type-parameters:1}

\subsection{Recall}\label{sec:02-functions-and-structures:02-type-parameters:2}

Formal definitions cited in this section, gathered here for quick reference (full citations in the \hyperref[sec:root:bibliography:1]{Bibliography}):

\begin{itemize}
\tightlist
\item
  \textbf{Type-parameterized structure.} This book's working statement: \texttt{structure Pair (α β : Type)} does not define one structure. It defines one structure \emph{per choice} of \texttt{α} and \texttt{β}, all at once. Categorically, this is a functor of the type parameters, not a single fixed structure.
\item
  \textbf{Parametric polymorphism.} Milner's theoretical account of type polymorphism (\cite{Milner1978}). Brief: genericity that is \emph{proved} once for every type, and checked before the generic code is ever called --- unlike an optional, erased type annotation. Python's \texttt{TypeVar} and Lean's \texttt{\{α : Type\} → ...} both implement this idea, to different degrees, as the Programmer's corner box below explains.
\end{itemize}

\begin{lstlisting}
structure Pair ((*$\alpha$*) (*$\beta$*) : Type) where
  fst : (*$\alpha$*)
  snd : (*$\beta$*)

def p : Pair Nat String := { fst := 1, snd := "one" }

#eval p.fst    -- 1
#eval p.snd     -- "one"
\end{lstlisting}

This generalizes directly to how we will write, e.g., \texttt{structure Group (α : Type)}.

\begin{mathreading}

\texttt{Pair α β} is the ordinary binary Cartesian product, but taken \emph{uniformly}, as a functor of two arguments at once. Rather than fixing \(\alpha\) and \(\beta\) once, \texttt{structure Pair (α β : Type) where ...} defines the \emph{whole family} of products \(\alpha \times \beta\) at once, one for every choice of the two type arguments. \texttt{Pair Nat String} is then just this construction evaluated at the pair \((\mathbb{N},
\mathrm{String})\), giving \(\mathbb{N} \times \mathrm{String}\). This is the same ``definition parameterized by objects of the category'' idea that lets us later write \texttt{Group (G : Type)}: not one group, but the functor sending a carrier type \(G\) to the type of group-structures on \(G\).

\end{mathreading}

\begin{progcorner}

Python code like \texttt{def identity(x): return x} or \texttt{class Pair: def \_\allowbreak{}\_\allowbreak{}init\_\allowbreak{}\_\allowbreak{}(self, fst, snd): ...} is also ``generic,'' in the sense that it does not mention any particular type. But nothing checks that genericity until the code actually runs. Writing \texttt{identity(3) + "oops"} causes Python to happily run \texttt{identity}, hand back \texttt{3}, and only then break on \texttt{3 + "oops"} at runtime. Constructing a \texttt{Pair} and placing a \texttt{Group (Fin 3)} into its \texttt{fst} draws no complaint from Python either.

\end{progcorner}

The closer Python analogue is \texttt{typing.TypeVar}: \texttt{def identity(x: T) -\allowbreak{}> T: return x} documents the same universally quantified type that \texttt{identity} has in Lean. But that documentation is optional and erased at runtime. It is enforced only if a checker such as \texttt{mypy} is separately run, and a stray \texttt{\# type: ignore} comment silences it entirely.

Lean's \texttt{\{α : Type\} → α → α} is not documentation. It is \emph{proved}, once, that \texttt{identity} works for every type \texttt{α}, and that proof is checked before \texttt{identity} is ever called --- not merely approximated by a linter that might never run.

\subsection{References}\label{sec:02-functions-and-structures:02-type-parameters:3}

Full citations in the \hyperref[sec:root:bibliography:1]{Bibliography}. Formal definitions are gathered in Recall, above.

\begin{itemize}
\tightlist
\item
  Python \texttt{typing} module documentation and mypy documentation (\cite{PythonTyping}, \cite{MypyDocs}) --- for the Python-side comparison used in this section's box.
\item
  Milner (\cite{Milner1978}) --- parametric polymorphism.
\end{itemize}

\section{Extending structures}\label{sec:02-functions-and-structures:03-extending-structures:1}

\begin{lstlisting}
structure Point3D extends Point where
  z : Nat

def origin3D : Point3D := { x := 0, y := 0, z := 0 }

#eval origin3D.x   -- inherited field, 0
\end{lstlisting}

\texttt{extends} is used later: a \texttt{CommGroup} (commutative group) will \texttt{extend} \texttt{Group} with one extra axiom (commutativity), instead of repeating all the group fields.

\begin{mathreading}

\texttt{structure Point3D extends Point where z : Nat} is the product \(\mathrm{Point3D} = \mathrm{Point} \times \mathbb{N} \cong
\mathbb{N} \times \mathbb{N} \times \mathbb{N}\), together with the \emph{forgetful map} \(\mathrm{Point3D} \to \mathrm{Point}\) (projecting away \(z\)), generated automatically as \texttt{.toPoint}. In algebraic language, this is exactly the pattern ``a \(\mathrm{Point3D}\)-structure is a \(\mathrm{Point}\)-structure plus one more piece of data.'' This is precisely how a \texttt{CommGroup} will later be ``a \texttt{Group}-structure plus one more axiom (\(ab = ba\))'': a full subcategory of \texttt{Group} cut out by an extra condition, with the forgetful functor \(\mathrm{CommGroup} \to \mathrm{Group}\) being exactly \texttt{.toGroup}.

\end{mathreading}

\subsection{References}\label{sec:02-functions-and-structures:03-extending-structures:2}

Full citations in the \hyperref[sec:root:bibliography:1]{Bibliography}.

\begin{itemize}
\tightlist
\item
  Lean 4 documentation (\cite{LeanDocs}) --- the auto-generated \texttt{.toX} projection produced by \texttt{extends}, used above as \texttt{origin3D.x} and \texttt{.toPoint}.
\item
  Mac Lane (\cite{MacLane1998}), Ch. I §3 ``Functors,'' p.~14 --- the categorical reading of the \texttt{.toPoint}/\texttt{.toGroup} projection generated by \texttt{extends}, verified verbatim: ``A functor which simply `forgets' some or all of the structure of an algebraic object is commonly called a forgetful functor (or, an underlying functor). Thus the forgetful functor \(U : \mathbf{Grp} \to \mathbf{Set}\) assigns to each group \(G\) the set \(UG\) of its elements\ldots{}''
\end{itemize}

\textbf{Key points.} A \texttt{structure} bundles data (and optionally proofs) under one name, built via the anonymous constructor \texttt{⟨...⟩} and read back out via field projection \texttt{.field}. \texttt{extends} builds a new structure containing everything an existing one has, plus more, generating a \texttt{.toX} forgetful projection for free --- the exact mechanism \texttt{CommGroup} uses to add commutativity to \texttt{Group} in Chapter 6.

\textbf{Socratic questions.}

\begin{enumerate}
\def\labelenumi{\arabic{enumi}.}
\tightlist
\item
  \emph{\texttt{⟨0, 0⟩} and \texttt{\{ x := 0, y := 0 \}} build the identical \texttt{Point}. Since the named form says more, why does the anonymous form ever get used?} Because the \emph{expected type} already says which fields are which --- \texttt{def origin : Point := ⟨0, 0⟩} cannot be ambiguous, since Lean already knows a \texttt{Point} is being built and in what field order. The named form earns its keep only when that context is not enough to make the assignment obvious to a reader.
\item
  \emph{\texttt{Point3D extends Point} generates \texttt{.toPoint} automatically. What would have to be written by hand instead, if \texttt{extends} did not exist?} A separate function \texttt{Point3D.toPoint : Point3D → Point} projecting out the shared fields one at a time --- exactly what a forgetful functor does explicitly, which is the whole reason \texttt{extends} is read categorically rather than as a mere convenience keyword.
\item
  \emph{A \texttt{structure} can bundle proofs as fields, not only data. What changes about }checking* that a term has the right type, once one of its fields is a proof rather than a number?* Nothing about the mechanism --- Lean still checks the field has the stated type --- but the stated type is now a proposition, so supplying that field means supplying a proof, checked once at construction time. This is exactly what makes \texttt{Group} (Chapter 6) impossible to build carelessly: the axiom fields cannot be filled in with nonsense that happens to type-check as data can.
\end{enumerate}


